
\color{red}
Współcześnie wynik ten ma ogólniejszą postać, zwaną twierdzeniem o jednorodności: w każdej płaszczyźnie Hilberta wszystkie trójkąty mają sumę kątów mniejszą od $\pi$, albo wszystkie równą $\pi$, albo wszystkie większą od $\pi$ (równoważnie: czwarty kąt każdego czworokąta Lamberta jest ostry, prosty albo rozwarty). % lang-pl
Według przypisu Hilberta, cytowanego przez Greenberga, udowodni to Max Dehn, a późniejsze dowody podadzą F.~Schur i Johannes Hjelmslev (zobacz sekcję \ref{sekcja_hjelmslev}); Saccheri i Lambert mieli ten wynik wcześniej, ale pierwszy z użyciem ciągłości \cite[s. 208]{greenberg_2010}. % lang-pl
W płaszczyznach typu ostrego i rozwartego trójkąty podobne są przystające, więc nie ma w nich teorii podobieństwa \cite[s. 208]{greenberg_2010}. % lang-pl
% Źródła: Hartshorne s. 326--333 (§36, zad. 36.1); por. Eves s. 301--304.
\color{black}
